\section{Prelocalized schobers on a disc}\label{section:prelocalized}
In this section, we study the naive categorification of $\Perv(\bbC,R)/\Loc(\bbC)$, which we denote by
\[\olsftwoP^{\pre}(\bbC,R) := \sftwoP(\bbC,R)/\sfLoc(\bbC).\]
As usual, we will make a choice of cuts and work at the level of abstract schobers.

In \S\ref{subsection:prelocalized}, we will define this quotient $\two$-category to be the right-orthogonal to the sub-$\two$-category of abstract local systems, $\sfSph^{\const}(n)\subseteq \sfSph(n)$.

In \S\ref{subsection:quotientismonad}, we show that there is an equivalence,
\[\sfSph/\sfSph^{\const}\simeq \sfSphMnd,\]
where $\sfSphMnd$ is the $\two$-category of pairs $(\Phi,M)$ consisting of a stable category $\Phi$ and a spherical monad $M$.

\begin{defn}\label{defn:sphmonad}
    A monad $M$ on category $\Phi$ is said to be a spherical monad if the monadic adjunction $S:\Phi\tofrom\Mod(M):R$ is a spherical adjunction.
\end{defn}

Spherical monads are studied in detail in \cite{christSphericalMonadicAdjunctions2023}. Observe that if $M$ is spherical, then $\Cone(\id\xto{u} M)$ is invertible. We will prove in Appendix~\ref{appendix:sphmonad}, that the converse is true, and so we have the following strengthening of \cite{christSphericalMonadicAdjunctions2023}*{Theorem 4.1}.

\begin{thm}\label{thm:christstrengthening1}
    A monad $M$ on $\Phi$ is spherical if and only if $\Cone(\id\xto{u} M)$ is invertible.
\end{thm}

For general $n$, we show that there is an equivalence
\[\sfSph(n)/\sfSph^{\const}(n)\simeq \sfSphMnd(n),\]
where $\sfSphMnd$ is the $\two$-category of pairs $(\frakS;\Phi,M)$ consisting of a stable category $\Phi$ and a spherical monad $M$ such that $\coCone(\id\xto{u} M)$ is conjugate-compatible with $\frakS$.

In \S\ref{subsection:twocatsphmnd}, we will upgrade the assignment
\[[n]\mapsto\sfSphMnd(n)\]
to a 2-Segal simplicial $\two$-category $\sfSphMnd(-)$, produce maps of simplicial $\two$-categories,
\[\sfSph(-)\to \sfSphMnd(-)\to \sfAut(-).\]
We then define a braid group action $\Br_n$ on $\sfSphMnd(n)$, compatibly with the braid group actions of $\sfSph(n)$ and $\sfAut(n)$.

\subsection{Orthogonal to local systems}\label{subsection:prelocalized}
We study the right-orthogonal to the sub-$\two$-category of local systems.  We begin by recalling which abstract perverse schobers are local systems.

\begin{defn}\label{defn:constant}
    An object $\calF\in\sfSph(n)$ is said to be an abstract local system if $\Phi_i=0$ for all $i$. For each $n$, we let $\sfSph^{\const}(n)\subseteq \sfSph(n)$ denote the full sub-$\two$-category consisting of abstract local systems.
\end{defn}

We study the quotient $\sfSph(n)/\sfSph^{\const}(n)$, which can be thought of as a naive categorification of abstract localized perverse sheaves. Since quotients of $\two$-categories are difficult, we will simply define the quotient to be the right orthogonal of $\sfSph^{\const}(n)$. We will later construct a projection functor as an adjoint to the embedding.

\begin{defn}
    An object $\calF\in\sfSph(n)$ is said to be an abstract prelocalized perverse schober with $n$ singularities if it is right orthogonal to $\sfSph^{\const}(n)$. We let $\sfSph^{\const\perp}(n)$ denote the full sub-$\two$-category of abstract prelocalized perverse schobers.
\end{defn}

\begin{prop}\label{prop:conservative}
    A perverse schober $(\Phi_i,\Psi,S_i,R_i)$ is an abstract prelocalized perverse schober if and only if the right adjoint
    \[\oplus_i\Phi_i\from\Psi:\oplus_iR_i\]
    is conservative.
\end{prop}

\begin{proof}
    We may identify objectwise, $\sfSt_k\isomto\sfSph^{\const}(n)$ via
    \[\opname{Const}:\calD\mapsto (\Phi_i=0,\Psi=\calD,S_i:0\to \Psi).\]
    Let $\calF:=(\Phi_i,\Psi,S_i:\Phi_i\to\Psi)$ be an abstract perverse schober. Observe that
    \[\Hom(\opname{Const}(\calD),\calF)\]
    is equivalent to the category of functors $G:\calD\to\Psi$ such that for each $i$, $R_i\circ G=0$.
    
    Suppose that $\calF$ is prelocalized. Let $\psi$ object of $\Psi$ such that $(\oplus_iR_i)(\psi)=0$. Then this defines an object of $\Hom(\opname{Const}(\Mod_k),\calF)$. We deduce that $\psi=0$, and hence that $(\oplus_iR_i)$ is conservative.

    Conversely, suppose that $(\oplus_iR_i)$ is conservative. Let $\calD$ be a category and let $f:\opname{Const}(\calD)\to\calF$ be a morphism. Given any object $d\in\calD$, $f(d)$ is an object of $\Psi$ satisfying
    \[(\oplus_iR_i)(f(d))=0.\]
    So by conservativity, $f(d)=0$. Hence $f=0$ and we deduce that $\calF$ is prelocalized.
\end{proof}

\begin{rmk}
    We compare this calculation to an analog with perverse sheaves. Let us take $k=\bbC$, $n=1$ for simplicity, and recall that $\Perv(\bbC,0)$ is equivalent to the category of diagrams of $k$-vector spaces,
    \[s:V\tofrom W:r\]
    such that $1-rs$ and $sr-1$ are invertible. We recall \cite{geiglePerpendicularCategoriesApplications1991} that for abelian categories, the correct notion of orthogonal subcategory asks for higher ext groups to vanish. With this in mind, an object $s:V\tofrom W:r$ is orthogonal to local systems if and only if $r$ is an isomorphism (and not just injective). This subcategory is easily seen to be equivalent to the category $\Loc(S^1)$.

    This is an instance in which the lack of a categorification of ``derived structure'' is a good thing; in the perverse schobers setting, asking for $R$ to be an equivalence (and not just conservative) would imply that $\Phi=\Psi=0$ and thus is not an interesting condition.
\end{rmk}

\begin{lem}\label{lem:activecartesian3}
    Let $\calF:=(\Phi_i,\Psi,S_i,R_i)$ be an object in $\sfSph(n)$. This is an abstract prelocalized perverse schober if and only if $\FS_{\sfSph}(\calF)$ is an abstract prelocalized perverse schober. Equivalently, there is a pullback square,
    \[\begin{tikzcd}
        {\sfSph^{\const\perp}(n)} \arrow[dr, phantom, "\scalebox{1.5}{$\lrcorner$}" , very near start, color=black] & {\sfSph^{\const\perp}} \\
        {\sfSph(n)} & {\sfSph.}
        \arrow[from=1-1, to=1-2]
        \arrow[hook, from=1-1, to=2-1]
        \arrow[hook, from=1-2, to=2-2]
        \arrow["{\FS_{\sfSph}}", from=2-1, to=2-2]
    \end{tikzcd}\]
\end{lem}

\begin{proof}
    Denote by $R$ the right adjoint appearing in $\FS_{\sfSph}(\calF)$. By Proposition~\ref{prop:conservative}, $\calF$ is a prelocalized schober if and only if $\oplus_iR_i$ is conservative, and $\FS_{\sfSph}(\calF)$ prelocalized if and only if $R$ is conservative. Observe that the following diagram commutes.
    % https://q.uiver.app/#q=WzAsMyxbMCwwLCJcXG9wbHVzX2lcXFBoaV9pIl0sWzEsMCwiXFxsYW5nbGVcXFBoaV8xLFxcZG90cyxcXFBoaV9uXFxyYW5nbGUiXSxbMiwwLCJcXFBzaSJdLFsxLDAsIkMiXSxbMiwxLCJSIl0sWzIsMCwiXFxvcGx1c19pUl9pIiwyLHsiY3VydmUiOjJ9XV0=
    \[\begin{tikzcd}
        {\oplus_i\Phi_i} & {\langle\Phi_1,\dots,\Phi_n\rangle} & \Psi
        \arrow["\oplus_i I_i^R", from=1-2, to=1-1]
        \arrow["{\oplus_iR_i}"', curve={height=12pt}, from=1-3, to=1-1]
        \arrow["R", from=1-3, to=1-2]
    \end{tikzcd}\]
    Here, $\oplus_i I_i^R$ is conservative by definition of SODs. It follows that $\oplus_iR_i$ is conservative if and only if $R$ is conservative, completing the proof.
\end{proof}

\subsection{Spherical monads}\label{subsection:quotientismonad}
We demonstrate a different description of $\sfSph^{\const\perp}$ in terms of spherical monads, and thereby exhibit an adjoint $\sfSph\to\sfSph^{\const\perp}$ to the inclusion of $\sfSph^{\const\perp}$.

We first ignore invertibility conditions and take $n=1$. The full sub-$\two$-category $\sfAdj^{\const\perp}\subset \sfAdj$ consists of those adjunctions where the right adjoint is conservative. These are equivalent to monadic adjunctions, so the embedding is identified with the Eilenberg--Moore construction,
\[\opname{P}^R:\sfMnd\inclto \sfAdj.\]
By definition, this is the right adjoint to the functor $P$, which outputs the monad associated to an adjunction. Note that in \S\ref{subsection:mndadj} we denoted these by $i^*\adjto i_*$.

Under this identification, $\sfSph^{\const\perp}$ is the sub-$\two$-category of those monads for which the image under $\opname{P}^R$ is spherical. Such monads are said to be spherical, as defined in Definition~\ref{defn:sphmonad}. We write $\sfSphMnd\subset \sfMnd$ for this sub-$\two$-category.

We deduce from Theorem~\ref{thm:christstrengthening1} that this restricts to an adjunction
\[\opname{P}:\sfSph\tofrom \sfSphMnd:\opname{P}^R.\]
There is also a functor $\sfSphMnd\to\sfAut$ given by $(\Phi,M)\mapsto (\Phi,\coCone(\id\to M))$.

For general $n$, we define
\[\sfSphMnd(n) := \sfAut(n)\times_{\sfAut}\sfSphMnd.\]
We denote its objects by $(\frakS;\Phi,M)$ consisting of a category $\Phi$, monad $M$, such that $\Cone(\id\to M)$ is conjugate-compatible with SOD $\frakS$, together with functors
\[\sfSph(n)\to\sfSphMnd(n)\to\sfAut(n).\]
We will show in Proposition~\ref{prop:monadisortho} that there is an identification,
\[\sfSph^{\const\perp}(n)\simeq \sfSphMnd(n).\]

\subsection{\texorpdfstring{$\sfSphMnd(n)$}{SphMnd(n)} as a simplicial \texorpdfstring{$\two$}{2}-category}\label{subsection:twocatsphmnd}
We construct the simplicial $\two$-category $\sfSphMnd(-)$. We will also construct a braid group action, and define the $\two$-category of non-abstract prelocalized schobers $\olsftwoP^{\pre}(\bbC,R)$. Since the techniques and ideas are the same as in previous sections, we will omit most details.

\begin{prop}\label{prop:simplicialsphmnd}
    The assignment $[n]\mapsto \sfSphMnd(n)$ upgrades to a simplicial $\two$-category which is 2-Segal. Moreover, the natural functors
    \[\sfSph(n)\xto{\opname{P}_n} \sfSphMnd(n)\to \sfAut(n)\]
    upgrade to maps of simplicial $\two$-categories which are Cartesian over $\Delta_{\act}$.
\end{prop}

\begin{proof}
    These follow in the same way as \S\ref{section:adjaut}. We outline the steps and omit the details.
    \begin{itemize}
        \item First consider the simplicial $\two$-category $\sfMnd'(-) := \sfEnd'(-)\times_{\sfEnd}\sfMnd$ and see that it is 2-Segal.
        \item Verify that the inclusion $\sfMnd\inclto \sfMnd'(1)$ admits a right adjoint via $(\Phi_1;\Phi,M)\mapsto (\Phi_1,I_1^RMI_1)$, by a proof similar to Lemma~\ref{lem:keyrightadjoint}, Lemma~\ref{lem:rightadjointend} and Lemma~\ref{lem:rightadjointbdj}.
        \item Verify that the inclusion $\sfMnd(n):= \sfEnd(n)\times_{\sfEnd}\sfMnd\inclto \sfMnd'(n)$ admits a right adjoint, by reducing to the case $n=1$.
        \item Use Lemma~\ref{lem:descendingsimplicialobjects} to upgrade $\sfMnd(-)$ to a simplicial $\two$-category. Verify that this is 2-Segal.
        \item See that the projection $\opname{P}:\sfAdj\to \sfMnd$ and the natural map $\sfMnd\to \sfEnd$ induce a map of simplicial $\two$-categories, $\sfBdj'(-)\to\sfMnd'(-)\to \sfEnd'(-)$.
        \item Use Lemma~\ref{lem:descendingsimplicialmaps} to obtain maps of simplicial objects $\sfAdj(-)\simeq\sfBdj(-)\to\sfMnd(-)$ and $\sfMnd(-)\to\sfEnd(-)$. Then verify that these are Cartesian over each active arrow $[1]\to [n]$ and hence over $\Delta_{\act}$.
        \item Finally, verify that these properties and structures all restrict to respective subcategories as given in the statement.
    \end{itemize}
\end{proof}

\begin{prop}\label{prop:monadisortho}
    For each $n$, the functor $\opname{P}_n$ admits a right adjoint $\opname{P}_n^R$ which is 2-fully faithful. Moreover, the resulting full sub-$\two$-category is $\sfSph^{\const\perp}(n)$.
\end{prop}

\begin{proof}
    Let $n=1$. We recall from the discussion at the beginning of this subsection that $\opname{P}:\sfSph\to\sfSphMnd$ admits a 2-fully faithful right adjoint $\opname{P}^R$ given by the inclusion of $\sfSph^{\const\perp}$.

    By applying $\sfAut(n)\times_{\sfAut}-$, we learn that $\opname{P}_n$ admits a 2-fully faithful right adjoint $\opname{P}_n^R$. Since $\sfSph(-)\to \sfSphMnd(-)$ is Cartesian over the active arrow $[1]\to [n]$, we have a pullback square as follows.
    \[\begin{tikzcd}[ampersand replacement=\&]
        {\sfSphMnd(n)}\pullback \& \sfSphMnd \\
        {\sfSph(n)} \& \sfSph
        \arrow[from=1-1, to=1-2]
        \arrow["{\opname{P}_n^R}"', hook, from=1-1, to=2-1]
        \arrow["{\opname{P}^R}"', hook, from=1-2, to=2-2]
        \arrow[from=2-1, to=2-2]
    \end{tikzcd}\]
    So we are done by comparing this to the diagram of Lemma~\ref{lem:activecartesian3}.
\end{proof}

From here on, we write $\sfSphMnd(n)$ instead of $\sfSph^{\const\perp}(n)$. As such, we refer to objects in $\sfSphMnd(n)$ as abstract prelocalized schobers and view $\opname{P}_n$ as the quotient functor. In Appendix~\ref{appendix:sphmonad}, we study the objects of $\sfSphMnd(n)$ in more detail.

\begin{cor}\label{cor:action3}
    There is a natural action of the braid group $\Br_n$ on $\sfSphMnd(n)$. Moreover, the two functors in Proposition~\ref{prop:simplicialsphmnd} are naturally $\Br_n$-equivariant.
\end{cor}

\begin{proof}
    Observe that the sub-$\two$-category $\sfSph^{\const}(n)$ is stable under the action of $\Br_n$ on $\sfSph(n)$. So we obtain an action of $\Br_n$ on $\sfSphMnd(n)$ which respects $\opname{P}_n^R$ and hence $\opname{P}_n$. It is easy to verify that this action also naturally respects the functor to $\sfAut(n)$. For example, we can factor this functor as $\sfSphMnd(n)\xto{\opname{P}_n^R}\sfSph(n)\to \sfAut(n)$.
\end{proof}

By the same construction as Definition~\ref{defn:nonabstract}, we obtain the $\two$-categories $\olsftwoP^{\pre}(\bbC,R)$ of prelocalized schobers, together with functors
\[\sftwoP(\bbC,R)\to \olsftwoP^{\pre}(\bbC,R)\to \olsftwoP(\bbC,R).\]